\chapter*{参考文献}
\addcontentsline{toc}{chapter}{参考文献}
\sourcepage{541--543}
文献代号、作者、期刊和出版信息按原书保留；题名后附中文译名。
\begin{description}[style=nextline]
\item[{[Ale24]}] James Waddell Alexander. \textit{An example of a simply connected surface bounding a region which is not simply connected}（围成非单连通区域的单连通曲面之例）. Proc. Nat. Acad. Sci. U.S.A., 10:8--10, 1924.
\item[{[Bro60]}] Morton Brown. \textit{A proof of the generalized Schoenflies theorem}（广义 Schoenflies 定理的证明）. Bull. Amer. Math. Soc., 66:74--76, 1960.
\item[{[Bro62]}] Morton Brown. \textit{Locally flat imbeddings of topological manifolds}（拓扑流形的局部平坦嵌入）. Ann. of Math. (2), 75:331--341, 1962.
\item[{[Can79]}] J. W. Cannon. \textit{Shrinking cell-like decompositions of manifolds. Codimension three}（流形类胞腔分解的收缩：余维三）. Ann. of Math. (2), 110(1):83--112, 1979.
\item[{[CF77]}] Richard H. Crowell、Ralph H. Fox. \textit{Introduction to knot theory}（纽结理论导引），Graduate Texts in Mathematics 57，Springer-Verlag，New York--Heidelberg，1977；1963 年原版重印。
\item[{[Dol95]}] Albrecht Dold. \textit{Lectures on algebraic topology}（代数拓扑讲义），Classics in Mathematics，Springer-Verlag，Berlin，1995；1972 年版重印。
\item[{[EM53]}] Samuel Eilenberg、Saunders MacLane. \textit{Acyclic models}（无环模型）. Amer. J. Math., 75:189--199, 1953.
\item[{[ES52]}] Samuel Eilenberg、Norman Steenrod. \textit{Foundations of algebraic topology}（代数拓扑基础），Princeton University Press，Princeton，1952。
\item[{[EZ53]}] Samuel Eilenberg、J. A. Zilber. \textit{On products of complexes}（论复形的乘积）. Amer. J. Math., 75:200--204, 1953.
\item[{[Fuc71]}] Martin Fuchs. \textit{A note on mapping cylinders}（关于映射柱的一则注记）. Michigan Math. J., 18:289--290, 1971.
\item[{[GH81]}] Marvin J. Greenberg、John R. Harper. \textit{Algebraic topology: A first course}（代数拓扑初步），Mathematics Lecture Note Series，CRC Press，Boca Raton，1981。
\item[{[Hil91]}] David Hilbert. \textit{Ueber die stetige Abbildung einer Linie auf ein Flächenstück}（论线到平面区域的连续映射）. Math. Ann., 38(3):459--460, 1891.
\item[{[Hop28]}] Heinz Hopf. \textit{Eine Verallgemeinerung der Euler--Poincaréschen Formel}（Euler--Poincaré 公式的推广）. Nachr. Akad. Wiss. Göttingen Math.-Phys. Kl. II, 127--136, 1928.
\item[{[Jor14]}] Camille Jordan. \textit{Cours d'analyse de l'école polytechnique. Tome III}（巴黎综合理工学院分析教程，第三卷），Cambridge Library Collection--Mathematics，Cambridge University Press，2014。内容为积分学与微分方程，1887 年第一版重印。
\item[{[Kne25]}] H. Kneser. \textit{Die Topologie der Mannigfaltigkeiten}（流形的拓扑）. Jahresber. Dtsch. Math.-Ver., 34:1--14, 1925.
\item[{[Man14]}] Ciprian Manolescu. \textit{Triangulations of manifolds}（流形的三角剖分）. ICCM Not., 2(2):21--23, 2014.
\item[{[Mas91]}] William S. Massey. \textit{A basic course in algebraic topology}（代数拓扑基础教程），Graduate Texts in Mathematics 127，Springer-Verlag，New York，1991。
\item[{[Maz59]}] Barry Mazur. \textit{On embeddings of spheres}（论球面的嵌入）. Bull. Amer. Math. Soc., 65:59--65, 1959.
\item[{[Maz61]}] B. C. Mazur. \textit{On embeddings of spheres}（论球面的嵌入）. Acta Math., 105:1--17, 1961.
\item[{[ML63]}] Saunders Mac Lane. \textit{Homology}（同调），Die Grundlehren der mathematischen Wissenschaften 114，Springer-Verlag，Berlin--Göttingen--Heidelberg；Academic Press，New York，1963。
\item[{[Moi52]}] Edwin E. Moise. \textit{Affine structures in 3-manifolds. V. The triangulation theorem and Hauptvermutung}（三维流形的仿射结构之五：三角剖分定理与主猜想）. Ann. of Math. (2), 56:96--114, 1952.
\item[{[Mor60]}] Marston Morse. \textit{A reduction of the Schoenflies extension problem}（Schoenflies 延拓问题的化简）. Bull. Amer. Math. Soc., 66:113--115, 1960.
\item[{[Mun00]}] James R. Munkres. \textit{Topology}（拓扑学），第二版，Prentice Hall，Upper Saddle River，2000。
\item[{[New51]}] M. H. A. Newman. \textit{Elements of the topology of plane sets of points}（平面点集拓扑基础），第二版，Cambridge University Press，1951。
\item[{[Pea90]}] G. Peano. \textit{Sur une courbe, qui remplit toute une aire plane}（论一条填满整个平面区域的曲线）. Math. Ann., 36(1):157--160, 1890.
\item[{[Poi10]}] Henri Poincaré. \textit{Papers on topology}（拓扑学论文集），History of Mathematics 37，American Mathematical Society，Providence；London Mathematical Society，London，2010。收录《位置分析》及五篇续篇，John Stillwell 翻译并撰写导言。
\item[{[Rei35]}] Kurt Reidemeister. \textit{Homotopieringe und Linsenräume}（同伦环与透镜空间）. Abh. Math. Sem. Univ. Hamburg, 11(1):102--109, 1935.
\item[{[Spa94]}] Edwin H. Spanier. \textit{Algebraic topology}（代数拓扑），Springer-Verlag，New York，1994；修订第三版重印。
\item[{[ST80]}] Herbert Seifert、William Threlfall. \textit{Seifert and Threlfall: a textbook of topology}（Seifert--Threlfall 拓扑学教材），Pure and Applied Mathematics 89，Academic Press，New York--London，1980。Joan S. Birman 作序，附 Seifert《三维纤维空间的拓扑》，Wolfgang Heil 从德文译出。
\item[{[Vic73]}] James W. Vick. \textit{Homology theory}（同调论：代数拓扑导引），Pure and Applied Mathematics 53，Academic Press，New York--London，1973。
\item[{[Vic94]}] James W. Vick. \textit{Homology theory}（同调论：代数拓扑导引），第二版，Graduate Texts in Mathematics 145，Springer-Verlag，New York，1994。
\item[{[Wey64]}] Hermann Weyl. \textit{The concept of a Riemann surface}（Riemann 曲面的概念），ADIWES International Series in Mathematics，Addison-Wesley，Reading--London，1964，德文版译本。
\item[{[Whi49]}] J. H. C. Whitehead. \textit{Combinatorial homotopy. I}（组合同伦，第一部分）. Bull. Amer. Math. Soc., 55:213--245, 1949.
\item[{[Whi78]}] George W. Whitehead. \textit{Elements of homotopy theory}（同伦理论基础），Graduate Texts in Mathematics 61，Springer-Verlag，New York--Berlin，1978。
\end{description}
